% !TeX program = xelatex
% !TeX encoding = UTF-8
% English editorial revision R20261107-02; authors and submission acceptance pending.
% Method first draft R20261106-06; Results R20261106-05; section plan R20261106-03.
% Standard natbib/BibTeX is used because local Biber cannot start; bibliography data remain in references.bib.
\documentclass[a4paper,11pt]{article}
\usepackage[english]{babel}
\usepackage{csquotes}
\usepackage{fontspec}
\usepackage[margin=2.5cm]{geometry}
\usepackage{amsmath,amssymb,graphicx,booktabs}
\usepackage{placeins}
\usepackage{xurl}
\usepackage[hidelinks]{hyperref}
\usepackage[numbers,sort&compress]{natbib}
\graphicspath{{figures/}}
\title{Home-Win Prediction in International Football:\\
Team Strength, Historical Information, and Regional Variation}
\author{Names, student IDs, and classes to be confirmed}
\date{}
\begin{document}
\maketitle
\begin{abstract}
Historical home-win rates in international football may reflect differences in team strength and match composition as well as venue-related conditions. We investigate how much these patterns contribute to pre-match prediction. A cohort of 24,917 non-neutral matches from 1980--2024 supports comparisons across eras, match types, and regions. A dynamic model learns strength from earlier results, including eligible neutral matches; a joint model then combines the strength distribution with match context. Retrospective chronological evaluation on 21,413 matches from 1990--2024 compares background-only, strength-only, and joint predictions. The joint model reduces log loss by 16.12\% relative to background alone and 0.42\% relative to strength alone, with 70.01\% binary accuracy for home-win/non-home-win decisions. Historical-window comparisons identify a trade-off between fitting cost and probability quality, without establishing a universally best length. Regional checks reveal uneven robustness and probability bias. A sealed release covers 52 future fixtures, with 49 main predictions and three withheld under the venue rule; outcomes remain unevaluated. Strength supplies most of the predictive gain, while context adds a modest increment whose interpretation is limited by historical reuse and model assumptions.
\end{abstract}
\noindent\textbf{Keywords:} international football; pre-match strength; home-win probability; probabilistic prediction
% !TeX root = ../main.tex
\section{Introduction}
% Introduce the motivation, related work, research questions, and contributions.
Draft section: substantive text has not been written.

% !TeX root = ../main.tex
\section{Method}
\label{sec:method}
% Editorial revision R20261107-02; implemented sources remain in method_evidence.json.

\subsection{Study Design and Data}
\label{sec:data-design}

We used the International Football Results dataset maintained by martj42~\cite{martj42_data}. The registered snapshot contains 49,945 matches dated 30 November 1872--6 October 2026, recording dates, teams, scores, tournaments, locations, and neutral-venue status. It is licensed under CC0 and identified by SHA-256; the upstream commit and original acquisition time are unknown. Scores include extra time where recorded and exclude penalty-shootout goals. The main outcome is $Y_i=1$ for a home win and $0$ for a draw or away win. A nominal home-team label at a neutral venue does not establish home advantage.

Preparation preserves the originals and records each row's disposition. Audits validate dates, integer scores, and venue flags, and flag duplicate candidates, conflicting scores, and reversed pairings. Critical unresolved records are quarantined; optional missing fields restrict eligibility for particular analyses. At the 31 December 2024 preparation cutoff, 48,102 matches were included, 1,803 were outside the window, and 40 were isolated. Identity and eligibility review then retained 24,917 of 26,055 study candidates, excluding 1,138. This includes 117 matches accepted following the project owner's manual verification, without a new independent official-source audit. Team identifiers denote football entities, which need not be sovereign countries; historical names and administrative changes are handled through explicit registry decisions.

The comparative cohort covers non-neutral matches in 1980--2024; earlier eligible records initialize strength. Eras are 1980--1989, 1990--1999, 2000--2009, 2010--2019, and the shorter 2020--2024 period. Normalized \texttt{Friendly} identifies friendlies; all other retained tournaments form the non-friendly group. Geographic labels begin with UN M49~\cite{un_m49}, with study overrides in both team roles: Turkey is in Europe/Southern Europe, Iran in Western Asia, and Australia in Asia/South-eastern Asia. Africa uses North, West, and remaining sub-Saharan groups; Oceania is not subdivided. These are fixed research labels, not historical confederation memberships.

Every match on a date is predicted before any result from that date updates strength, because reliable intraday order is unavailable. Outcome-model coefficients are refitted annually using previous years only. Strength learning uses 38,919 eligible matches from 1950: 5,581 pre-1980 matches, 8,421 later neutral matches, and the 24,917-match comparative cohort. Initial development and selection used 1970--1974 and 1975--1979, respectively, but later diagnostics and revisions reused subsequent history. Evaluation is therefore retrospective. Figure~\ref{fig:study-design} shows the overlapping learning and comparison samples.

\begin{figure}[htbp]
\centering
% Native vector diagram for the Method first draft, R20261106-06.
% Counts denote distinct, overlapping uses; arrows are information flow, not subtraction.
\setlength{\unitlength}{1mm}
\begin{picture}(150,61)
\put(0,50){\framebox(150,11){\shortstack{\small Registered results snapshot: 49,945 matches\\\small Deterministic preparation and explicit identity/eligibility rules}}}
\put(50,50){\vector(-1,-1){7}}
\put(100,50){\vector(1,-1){7}}
\put(0,28){\framebox(70,15){\shortstack{\small Comparative cohort: 24,917\\\small Non-neutral, 1980--2024\\\small Historical rates and fixed references}}}
\put(80,28){\framebox(70,15){\shortstack{\small Strength learning: 38,919\\\small Eligible matches, 1950--2024\\\small Includes earlier and neutral matches}}}
\put(35,28){\vector(0,-1){9}}
\put(115,28){\vector(0,-1){9}}
\put(0,2){\framebox(150,17){\shortstack{\small Historical comparisons $\longrightarrow$ strength-informed probability analysis\\\small Time-forward evaluation $\longrightarrow$ history/region checks $\longrightarrow$ prospective release\\\small Distinct targets and samples; see Table~\ref{tab:analysis-samples}}}}
\end{picture}

\papercaption{Information flow and analytical branches.}{The 24,917-match comparative cohort and the 38,919-match strength-learning set overlap; neutral and pre-1980 matches contribute to learning without entering the principal historical home-win comparisons. Only actual earlier results update later states. The sealed prospective release uses its own stated cutoff.}
\label{fig:study-design}
\end{figure}

\subsection{Historical Comparison Framework}
\label{sec:comparison-method}

For a group $g$ with $n_g$ observed matches, the direct rate is $\widehat r_g=n_g^{-1}\sum_{i\in g}Y_i$. To compare different compositions, define strata $c$ and a common reference distribution $w_c\geq0$, $\sum_c w_c=1$. We calculate
\begin{equation}
 \widehat r_g^{\,\mathrm{ref}}=\sum_c w_c\widehat r_{gc},\qquad
 \widehat\Delta_{g,h}=\widehat r_g^{\,\mathrm{ref}}-\widehat r_h^{\,\mathrm{ref}}.
 \label{eq:standardization}
\end{equation}
Era comparisons fix the friendly/non-friendly composition; match-type comparisons fix the era composition. Regional comparisons use their specified era/type reference and retain the direction from home group to opponent group. A missing, positively weighted cell makes a comparison unsupported. Ten-year windows advance one calendar year and overlap. The descriptive sensitivity analysis of friendly weights is separate from their information weight in strength learning.

Observed proportions describe a fixed history. Uncertainty assessment concerns the average tendencies that could generate outcomes under the same schedule, groups, denominators, and reference. Each comparison is a weighted sum $\widehat\theta_j=a_j^\top Y$. Shared teams and repeated pairings motivate dependence modelling~\cite{aronow2015dyadic}; our procedure uses a specified covariance model rather than a conventional dyadic sandwich estimator.

Mean fitting removes part of the variation used to estimate dependence. To account for this, let $X$ contain the fixed mean-model terms, $H=X(X^\top X)^+X^\top$, and $M=I-H$. Under $E[Y]=X\beta$, residuals satisfy $E[(MY)(MY)^\top]=M\Sigma M^\top$. We therefore fit covariance components after the same residual projection. Six non-negative components represent independent noise, repeated pairs, same/opposite-role team processes, and two shared calendar-time processes. A scale $s_j(Y)$ combines estimated quadratic variation with predefined protection against slow common changes. Separate complete-history simulations calibrate a maximum-error multiplier $c_F$ within each prespecified target family, giving
\begin{equation}
 I_j=\bigl[\widehat\theta_j-c_Fs_j(Y),\widehat\theta_j+c_Fs_j(Y)\bigr]\cap\mathcal L_j,
 \label{eq:conditional-envelope}
\end{equation}
Here $\mathcal L_j$ is the logical range and $s_j$ is not a conventional standard error. Final rank calibration uses error parameter $0.01$ separately in four prespecified families of 27, 64, 36, and 71 quantities. Its guarantee concerns frozen reference generators over calibration and new synthetic histories, not all 198 quantities jointly or arbitrary football processes. Detailed assumptions and numerical validation gates are indexed in the supporting-results guide.

The uncertainty design addresses a specific failure of treating matches as independent Bernoulli trials. Even with the same marginal win probabilities, $\operatorname{Var}(a_j^\top Y)=a_j^\top\Sigma a_j$ includes off-diagonal terms from shared teams and calendar conditions. Large match counts do not force these terms to vanish. The residual projection above also shows why fitting covariance directly to residual products as if their expectation were $\Sigma$ can understate variation: their expectation is $M\Sigma M^\top$. Our component model corrects that fitting target, but still restricts which dependence structures can occur. Sparse comparison cells cannot be repaired by a more elaborate covariance estimator.

We compared covariance estimation and calibration routes against common simulated reference truths. The adopted historical method estimates both mean and long-term common variation, then calibrates family maxima over complete generated histories. This choice protects the specified simultaneous comparisons more directly than reporting independent binomial intervals or an uncalibrated residual-only scale. Its cost is sometimes a very broad range, especially for era changes observed along only one historical path. A permanent global-shock generator remains a documented failure boundary: its recorded conditional coverage was 866 of 1,000 histories. We therefore restrict claims to admitted mechanisms rather than interpreting successful calibration as a distribution-free guarantee. Likewise, overlapping ten-year windows visualize continuity, but do not constitute independent repetitions that increase the effective sample size tenfold.

% !TeX root = ../main.tex
\subsection{Model Choice and Mathematical Rationale}
\label{sec:model-choice}

The model must learn from sparse, irregularly spaced results, update both opponents, distinguish draws, and represent uncertainty before a result is known. It must also separate team ability from the advantage of playing at home. These requirements determine which parts of earlier work we use. Elo-derived football predictors provide a useful separation between ratings and outcome regression~\cite{hvattum2010elo}, but a fixed-step point rating does not itself describe uncertainty about ability. We retain a fixed-step internal baseline for comparison, rather than interpreting the numerical difference between two rating points as an observed physical difference in football quality. Our latent difference is a coordinate in a probability model; its unit is fixed by performance-noise variance, and its effect on probability is nonlinear.

Dynamic paired comparisons~\cite{glickman1999dynamic} supply the central distinction between latent ability and uncertainty about that ability. A Gaussian evolution model is a tractable starting point: inactivity enlarges uncertainty while leaving the expected strength unchanged. However, a single smooth diffusion may adapt poorly to persistent abrupt changes. We therefore develop both the smooth special case and a model with an additional persistent jump component. The jump model is selected for the admitted development requirements, not because mathematical complexity guarantees better prediction. A simpler model can remain preferable in a stationary simulated setting.

TrueSkill~\cite{herbrich2006trueskill} offers related ideas of uncertain skill, noisy performance, and a draw region. These ideas suit a result archive containing home wins, draws, and away wins. Its game-rating application does not by itself establish validity for national teams. Our dated evolution, venue term, robust likelihood, and chronological checks are additional project decisions. Expectation propagation~\cite{minka2001ep} supplies an approximate way to communicate outcome information through correlated latent states; it is a computational method, not a guarantee that the entire posterior is Gaussian or exact.

A goal-count model, illustrated by Dixon and Coles~\cite{dixon1997scores}, could use more of the recorded score information. We do not adopt it for this analysis: the principal outcome is a home win, and scores are not uniformly restricted to regulation time. Adding a score model would require its own observation assumptions and evaluation rather than treating goals and the resulting win as independent evidence. FIFA rankings and published Elo ratings are not tested external benchmarks here. Neither the present comparisons nor the cited papers establish that our implementation outperforms those systems.

Two other method choices concern different problems. Direct standardization answers how observed rates differ under a common match composition; it does not learn team strength. Proper scoring rules~\cite{gneiting2007scoring} evaluate whether probabilities are useful for subsequent outcomes, while simulation-based uncertainty analysis examines stated comparison procedures under specified generating mechanisms. No one of these tools substitutes for the others. In particular, model-conditional interval coverage cannot establish real-world predictive superiority.


\subsection{Dynamic Pre-Match Team Strength}
\label{sec:strength-method}

Strength is an unobserved relative state inferred from results. Dynamic paired comparisons distinguish changes in ability from uncertainty about it~\cite{glickman1999dynamic}. We combine gradual Gaussian changes with persistent Poisson jumps. For team $u$, over elapsed time $\Delta$ in years (days divided by $365.2425$),
\begin{equation}
 s_u(t+\Delta)=s_u(t)+G_u+\sum_{\ell=1}^{N_u}J_{u\ell},\quad
 G_u\sim N(0,q\Delta),\quad N_u\sim\mathrm{Poisson}(\lambda\Delta),\quad
 J_{u\ell}\sim N(0,\tau^2).
 \label{eq:strength-transition}
\end{equation}
The innovations have zero prior mean. Inactivity therefore increases uncertainty without mechanically lowering the strength centre. A jump is a persistent change in state, not an additional label assigned only to an upset.

The two components encode different hypotheses about change, not two types of observed match. Independence and the law of total variance give
\begin{equation}
 E[s_u(t+\Delta)-s_u(t)]=0,\qquad
 \operatorname{Var}[s_u(t+\Delta)-s_u(t)]=(q+\lambda\tau^2)\Delta.
 \label{eq:transition-moments}
\end{equation}
Indeed, conditional on $N_u=n$, the jump sum has mean zero and variance $n\tau^2$; averaging over $E[N_u]=\lambda\Delta$ gives the second term. Matching this variance with a single Gaussian would not preserve the jump mixture's tails. A long absence therefore permits a wider set of strengths without supplying evidence that the team has improved or declined. Once play resumes, actual results determine the direction of learning.

For a match between home team $h$ and away team $a$, the performance variable is $Z_i=s_h-s_a+v_iH_t+\varepsilon_i$, with $\varepsilon_i\sim N(0,1)$, shared home term $H_t$, and venue sign $v_i\in\{-1,0,1\}$. Neutral venues have $v_i=0$. Outcomes correspond to $Z_i>\delta$ (home win), $|Z_i|\leq\delta$ (draw), and $Z_i<-\delta$ (away win). We mix these informative probabilities with a uniform three-outcome component:
\begin{equation}
 P(R_i=r\mid\mathrm{state})=\kappa_i P(Z_i\in A_r\mid\mathrm{state})+(1-\kappa_i)/3,\qquad
 \kappa_i=(1-\epsilon)\begin{cases}\alpha,&\text{friendly},\\1,&\text{otherwise}.\end{cases}
 \label{eq:rating-likelihood}
\end{equation}
The mixture reduces how informative a friendly result is about strength. Its coefficient does not directly multiply the win rate, match count, or rating update. The same likelihood handles every result: a single surprise can reflect noise, while repeated evidence can support a persistent change.

For a given latent state, write $d=s_h-s_a+vH$. The informative probabilities are $\ell_H=\Phi(d-\delta)$, $\ell_A=\Phi(-\delta-d)$, and $\ell_D=\Phi(\delta-d)-\Phi(-\delta-d)$. The difference is thus an input to a nonlinear probability map, not a claim that a one-unit ability advantage always produces the same win-rate change. Integrating these probabilities over the uncertain state also distinguishes a well-supported advantage from an equally large but poorly known one.

Bayes' rule makes the handling of surprises explicit. If $S$ denotes the full pre-match state, $f^-(S)$ its prior density, and $r$ the observed result, then
\begin{equation}
 f^+(S)=\frac{f^-(S)\{\kappa\ell_r(S)+(1-\kappa)/3\}}{\int f^-(u)\{\kappa\ell_r(u)+(1-\kappa)/3\}\,du},\qquad
 \nabla\log P(r\mid S)=w_r(S)\nabla\log\ell_r(S),
 \label{eq:robust-learning}
\end{equation}
where $w_r(S)=\kappa\ell_r(S)/\{\kappa\ell_r(S)+(1-\kappa)/3\}$. This is the posterior responsibility of the informative component. An outcome very unlikely under a particular state has small responsibility there, reducing its leverage relative to the informative likelihood alone. The multiplier is bounded between zero and one, but the update is not a fixed linear multiple of the match result: responsibilities, state probabilities, and posterior moments change together. Other states, including jump-supported states, may explain the result better. Repeated surprising results can accumulate support for a lasting shift, although no theorem here guarantees recovery within a specified number of games. Stronger robustness can also delay recognition of a genuine decline; this trade-off requires diagnostics rather than an upset-specific exception.

The adopted parameters are $q=0.03$, $\lambda=0.08$ per year, $\tau=1.2$, $\delta=0.4$, $\alpha=0.65$, and $\epsilon=0.05$; the shared home term has variance rate $0.005$ per year. Team priors have mean zero and variance one; the home-term prior has mean zero and variance $0.25$. Unit performance-noise variance defines the rating scale. Parameter and initialization choices were assessed through finite development checks.

These values are empirical design choices, not constants deduced from the cited rating papers. The initial global protocol registered 21 settings, including friendly information coefficients $0.4$, $0.65$, and $0.85$, and compared loss, class bias, structural complexity, and numerical admissibility. The jump challenger initially required further numerical and cost work before admission; its later repair and probability-correction selection reused history. Thus $\alpha=0.65$ has a finite model-development basis, but it is neither a measured universal fraction of friendly-match effort nor a globally optimal weight. The final eight-way correction comparison used a separate frozen candidate list and unchanged acceptance gates. The development sequence, including its failures, is reported in Section~\ref{sec:strength-development} rather than recast as a single successful pre-1980 selection.

These parameters identify the adopted full-jump rating input. Its authoritative parameter record is the registered rating report, distinguished from the zero-jump special case used in an earlier handoff rule-evaluation exercise. Table~\ref{tab:predictive-models} uses the full-jump input and the v2 annual prediction protocol; the supporting guide gives the exact configuration and output identifiers.

At each date, Bayesian updating combines outcome likelihoods with the joint prior, retaining means and covariances for both opponents. Poisson-count likelihoods use controlled tail truncation; a team's matches on the same date share its latent changes. Damped expectation propagation and end-of-date Gaussian projection provide an approximate sequential calculation~\cite{minka2001ep}. Both teams learn from each result without a separate away-strength parameter.

The covariance is essential: for opponents with marginal variances $V_h,V_a$ and covariance $C_{ha}$, $\operatorname{Var}(s_h-s_a)=V_h+V_a-2C_{ha}$. Treating them as independent would generally change predictive uncertainty. At a single Gaussian-prior update, the moment identities $m^+=m+V\nabla_m\log\mathcal Z$ and $V^+=V+V(\nabla_m^2\log\mathcal Z)V$ follow by differentiating the integrated likelihood $\mathcal Z$ with respect to the prior mean. They explain why likelihood evidence updates both the centre and uncertainty. They are not a claim that repeated mixture updates remain exactly Gaussian. The implementation first integrates local jump counts, jointly treats matches sharing a team on the date, then projects the resulting joint moments at the day boundary. This avoids expanding every team's jump history as a global collection of branches. The unresolved loss of multimodality and the lack of a general EP convergence guarantee remain approximation limits.

We initialize in 1950 rather than impose an arbitrary downweighting on all pre-war results. This is an implemented scope choice, not a mathematically optimal historical boundary. The 5,581 eligible 1950--1979 games give the model a learning period before the principal cohort; dynamic changes reduce the need to treat old outcomes as equally informative about current ability. Starting in 1960 and changing the initial variance to $0.5$ or $2$ were also checked. Cold starts and historically weak links between opponent groups remain flagged rather than being removed or declared solved by a common prior.

\subsubsection{Online Learning and the Separate Probability Correction}

This is a form of supervised probabilistic learning: an observed result is the label, and posterior states are updated for subsequent predictions. The latent-strength model is not a neural network. Its finite parameter search and the outcome layer's numerical optimization are different operations from updating a team's state after a match. Predicting first and learning later prevents the current result from explaining its own prediction; it does not remove bias from reusing history to choose a model.

Development exposed probability-shape errors and persistent recent residual bias, which require different corrections. For the auxiliary three-class branch, let $p_H,p_D,p_A$ be raw rating probabilities and $u=(\log(p_H/p_A),\log(p_D/p_A))$. Within the registered friendly/neutral contexts, the adopted annual correction uses the preceding ten years to fit a shared slope $a$ and two offsets $c$. For fixed fitted parameters,
\begin{equation}
 (\widetilde p_H,\widetilde p_D,\widetilde p_A)
 =\operatorname{softmax}(a u_H+c_H,\ a u_D+c_D,\ 0).
 \label{eq:shape-correction}
\end{equation}
The regularized likelihood shrinks toward the identity map $(a,c)=(1,0)$, with sparse contexts falling back to their registered parent group. Parameter uncertainty is handled approximately by local Gaussian integration. A second, zero-drift two-offset state tracks remaining time-local bias, with diffusion $0.02$ per year and same-day joint updates. Its predictive probabilities integrate the offsets before learning the day's outcomes. These corrections change probabilities, not the underlying team means or covariances. They must not be fed into the joint model as another supposedly independent strength observation. The joint model below instead receives the uncorrected latent team-difference distribution, excludes the rating model's home term, and estimates its own context effects.

\subsection{Joint Probability Model}
\label{sec:joint-method}

The joint model receives the \emph{team-strength difference} distribution, with the filter's shared home term excluded so that context is handled in the outcome model. Write the pre-match difference as
\begin{equation}
 D_i\mid K_i=k\sim N(\mu_i,\sigma_{D,i}^2+k\tau^2),\qquad
 K_i\sim\mathrm{Poisson}(\nu_i),
 \label{eq:difference-mixture}
\end{equation}
where $\sigma_{D,i}^2$ is the base variance of the difference, including the opponents' covariance, and $\nu_i$ is their expected total count of unobserved jumps. With Poisson weights $\pi_k(\nu)=e^{-\nu}\nu^k/k!$, integrating the probit predictor over uncertain strength gives
\begin{equation}
 p_i=\sum_{k=0}^{\infty}\pi_k(\nu_i)
 \Phi\!\left(\frac{\eta_i+b\mu_i}{\sqrt{1+b^2(\sigma_{D,i}^2+k\tau^2)}}\right),
 \qquad \eta_i=x_i^\top\gamma,\quad b>0.
 \label{eq:joint-probability}
\end{equation}
Here $\Phi$ is the standard normal distribution function. Context $x_i$ includes era, friendly status, both teams' geographic regions, and within-region groups centred by their frequencies. The estimated positive slope $b$ maps strength to the probability scale; integrating over the distribution avoids treating estimated strength as known. Background-only uses context without strength ($b=0$); strength-only retains an intercept and strength while excluding other contextual terms.

The denominator in Equation~\ref{eq:joint-probability} follows analytically rather than from a fitted uncertainty penalty. Conditional on a jump count, let $D\sim N(\mu,v)$ and independently $U\sim N(0,1)$. Then
\begin{equation}
 E_D[\Phi(\eta+bD)]
 =P\{U-b(D-\mu)\leq\eta+b\mu\}
 =\Phi\!\left(\frac{\eta+b\mu}{\sqrt{1+b^2v}}\right).
 \label{eq:probit-integral}
\end{equation}
Applying this identity at $v=\sigma_D^2+k\tau^2$ and averaging over $k$ gives the adopted mixture. Larger uncertain variance moderates extreme probabilities conditional on the same numerator; it is not a rule that every uncertain match becomes a draw. This binary layer does not separately assign draw and away-win probabilities.

For two allowed context scenarios $a,b$, their strength-adjusted contrast is
\begin{equation}
 \widehat\Delta^{\mathrm{str}}_{a,b}
 =\widehat E_{\mathcal R}\!\left[\widehat p(X^{(a)},D)-\widehat p(X^{(b)},D)\right].
 \label{eq:strength-contrast}
\end{equation}
Here $\widehat E_{\mathcal R}$ uses the registered reference weights and scenario mixtures while retaining the reference strength distributions. Computing the corresponding contrast before and after adding strength gives the paired change reported in the results. This expresses a model-based comparison, not a change in physical venue or team identity.

Two fitting tasks are distinct. Conditional analysis fits the 24,917-match cohort and averages paired predictions under allowed context combinations on a fixed reference. Historical prediction refits using matches from 1980 through the preceding year and evaluates 1990--2024. Geographic coding uses that year's training rows only; unseen within-region groups receive the region-centred zero effect. The annual objective is mean Bernoulli negative log likelihood plus $\|\theta\|^2/(2n)$, where $\theta$ includes contextual coefficients and log slope; conditional analysis is unpenalized. L-BFGS-B uses initial slopes $0.3,1,2$, coefficient bounds $[-8,8]$, and slope bounds $[10^{-4},8]$, with gradient and multistart checks. Poisson tails use tolerance $10^{-12}$, checked separately at tighter tolerance.

Fixed-reference comparisons retain flags for empty or sparse cells and out-of-range strength. Narrower overlap references are sensitivity analyses. Auxiliary win--draw--loss forecasts are a separate branch: temperature/offset fitting over the past ten years and an online residual correction act on native probabilities, with residual diffusion $0.02$ per year. They do not alter the joint model's strength input or its binary target.

\subsection{Evaluation, Uncertainty, and Implementation}
\label{sec:evaluation-method}

For the same evaluation matches, we measure probability quality by mean log loss $L$ and Brier score $B$:
\begin{equation}
 L=-\frac1n\sum_i\{Y_i\log p_i+(1-Y_i)\log(1-p_i)\},\qquad
 B=\frac1n\sum_i(p_i-Y_i)^2.
 \label{eq:probability-scores}
\end{equation}
Lower scores indicate better probability forecasts~\cite{gneiting2007scoring}. We also report binary accuracy with home win chosen at $p_i\geq0.5$, and bias $100(\bar p-\bar Y)$ in percentage points (pp). Relative loss reduction is $(L_{\mathrm{reference}}-L_{\mathrm{joint}})/L_{\mathrm{reference}}$, distinct from an accuracy-point gain. Loss comparisons use paired matches; no real-data dependence-adjusted interval for the mean difference is available.

Three uncertainty statements serve different purposes: the rating distribution represents uncertain strength; simulation assesses conditional comparisons; match probabilities describe uncertain outcomes. The prospective central 90\% ranges propagate latent strength uncertainty with coefficients fixed, excluding coefficient and model error.

Conditional contrasts require additional validation because outcomes update the strength estimates used later. Route A generates a full history, updates strength, and refits the analysis; route B uses a first-order total-score approximation. Separate reference-truth, calibration, and validation streams assess maximum standardized error across the comparisons, allowing for Monte Carlo truth error. We adopt A and report only quantities meeting the simulation acceptance criteria. Generators include date-conditionally independent outcomes and a same-day latent common-shock setting of $0.25$, which is a generator parameter rather than the correlation of all match pairs. Follow-up simulations use a preallocated sequential error budget, following time-uniform control principles~\cite{howard2021confidence}. This validation concerns conditional contrasts under those generators, separately from predictive loss comparisons.

The two routes account for a feedback absent from a regression with externally known strength: changing an early result can change later ratings, later predictions, and fitted context coefficients. Route A recomputes that chain on each generated history. Route B approximates the change through a first-order expansion of the full score, Hessian, and target gradient; its validity also depends on small linearization remainder. A preliminary date-multiplier approach did not establish control of predictable mean drift and long-term feedback, so it was not accepted as the final B route. Both final routes used the same generators and truth targets. They showed close numerical agreement, but A was adopted because complete refitting avoids substituting a first-order approximation for the full change. B's optimistic total-cost saving in two representative histories was about $8.93\%$, below the prespecified $20\%$ replacement threshold. This is a methodological reason supported by finite comparisons, not proof that A dominates every alternative or that two agreeing routes give independent real-world evidence. Finite Monte Carlo truth error still prevented admission of some regional changes, as reported in the results.

History-window comparisons vary strength memory and outcome-learning history separately. Finite-memory states are reconstructed annually; an adaptive choice uses past performance only. Outcome parameters stay fixed within each evaluation year while strength updates from earlier actual results. Timings distinguish the measured outcome-fitting component from total costs. Regional diagnostics hold predictions fixed, standardize nine cells combining strength-gap and strength-uncertainty categories, and use cutpoints learned before 2010. Leave-one-year/team checks omit affected evaluation matches, in either team role, without rebuilding the learning history.

The adopted prediction manifest records Python 3.14.8, NumPy 2.5.4, and SciPy 1.18.1 on Windows 11. Exact recorded dependencies and stage-specific random-stream settings are indexed in the supporting guide; this is an execution record, not a claim of reproduction on another machine. The Python implementation uses NumPy for arrays and covariances, SciPy for distributions and optimization, and Matplotlib for figures~\cite{numpy2020,scipy2020,matplotlib2007}. Repository records identify adopted inputs, configurations, source hashes, environment, and stage entry points. Evidence registers link manuscript claims to those records; detailed settings, identity/group mappings, derivations, and failed candidates are retained for inspection and reproduction.

% !TeX root = ../main.tex
\section{Experimental Results and Analysis}
% Organize the evidence by research question; identify the sample and scope.
Draft section: substantive text has not been written.

% !TeX root = ../main.tex
\section{Discussion and Conclusion}
% Discuss findings, contributions, limitations, lessons, and proposed extensions.
Draft section: substantive text has not been written.

\clearpage
\begingroup
\setlength{\emergencystretch}{2em}
\setlength{\bibsep}{2pt}
\bibliographystyle{unsrtnat}
\bibliography{references}
\endgroup
\end{document}
